\section{The stationary state}
\label{sec:steady}

Marching to a steady state is only one of the two ways EXHALE reaches one.
The other is to solve for it directly, as the root of the stationary residual
of the very same discrete operator.  This section states that residual, the
measure by which a state is judged steady, and the two solvers that drive it
to zero.

\subsection{The stationary residual}
\label{sec:steady:residual}

\texttt{assemble\_residual} (\texttt{steady\_residual.f90}) evaluates the
right-hand side of one forward Euler step of the marching operator and
subtracts the radiative source from its third row,
%
\begin{equation}
  \begin{aligned}
    R_1 &= dF_1 - S_1, \\
    R_2 &= dF_2 - S_2 - S_\mu, \\
    R_3 &= dF_3 - S_3 - (\mathcal{H} - \Lambda) - S_{\mathcal{E}}
           + \nabla\!\cdot q_d,
  \end{aligned}
  \label{eq:residual}
\end{equation}
%
with $\nabla\!\cdot q_d$ the divergence of the interdiffusion enthalpy flux of
Sec.~\ref{sec:mol:interdiffusion}, formed from the element fluxes of the
state itself and identically zero, with no arithmetic, unless
\texttt{He\_diffusion} moves the elements; the composition is therefore an
argument of every assembly.  $dF$ and $S$ are those of \eqref{eq:dF1} to
\eqref{eq:source}, evaluated through
the same \texttt{reconstruction\_continuation\_rhs} the marching stages call,
so that the fixed point of the marching and the zero of the residual are the
same state by construction.  The lowest ghost cell has no lower face and
therefore no balance to state; its row is explicitly zero, so that no reader
of $R$ finds a heating rate standing alone in a row that has no fluxes.

A residual row is a dimensional quantity and has to be divided by something
before it can be compared with a tolerance.  Every row is divided by the
largest term \emph{that row itself contains}
(\texttt{residual\_row\_scale}):
%
\begin{equation}
  \begin{gathered}
  s_1(j) = \frac{\max\bigl(|F_{j-1/2}|A_-,\ |F_{j+1/2}|A_+\bigr)}{\Delta V_j},
  \quad
  s_2(j) = \max\bigl(|{\rm ram}|,\ |\partial_r p|,\
                     \rho|\partial_r\Phi|,\ |S_\mu|\bigr),
  \\
  s_3(j) = \max\bigl(|dF_3|,\ |S_3|,\ \mathcal{H},\ \Lambda,\
                     |S_{\mathcal{E}}|,\ |\nabla\!\cdot q_d|\bigr).
  \end{gathered}
  \label{eq:rowscales}
\end{equation}
%
Since $\Delta V_j$ is exactly the volume the flux difference was divided by,
$R_1/s_1$ is the fractional change of the mass flux across the cell, and the
mass row having no source, that fraction is the whole of what the row says.
The momentum scale uses the three gathered physical terms of
Sec.~\ref{sec:hydro:pressure} and not $dF_2$ and $S_2$, for the reason given
there.

A signal-speed bound $D_k(|v|+c_s)/\Delta r$ is not usable as the scale of
these rows: nothing in any of them moves at $c_s$, so in a quasi-hydrostatic
layer that bound stands far above the terms the row holds and $R_1$ divided by
it is the fractional flux error per cell \emph{times} the local Mach number,
which at Mach $5\times10^{-5}$ is a blindness of four decades.

These are the scales of the \emph{measure} only.  The Newton system and its
merit keep the state scales $D_k$ of \texttt{state\_scales\_of\_cell}, so the
merit and the acceptance test are deliberately not one expression apart.

\subsection{The acceptance test}
\label{sec:steady:gates}

\texttt{steady\_gates\_met} is the only definition of an accepted steady state
in the code: every route that may stop on ``the equations are satisfied'', the
two stationary solvers and the marching loop under \texttt{Resid tol:}, calls
it.  Four gates:
%
\begin{enumerate}
\item the \emph{residual gate}, $\|R\| < $ \texttt{resid\_tol}, with $\|R\|$
      the maximum over physical cells and rows of $|R_k(j)|/s_k(j)$;
\item the \emph{flux gate}, the spread of the Riemann face mass flux
      $F_{f+1/2}r_{f+1/2}^2$ over $r \ge r_{\rm flux}$ below
      \texttt{flux\_spread\_th} (default $2\times10^{-5}$ at
      $r_{\rm flux} = 1.2\,R_p$).  A residual norm does not say this: a
      slowly varying flux has a small local derivative, and two states that
      differ by a factor 25 in the spread differ by only 1.9 to 5.7 in every
      residual norm tried.  It is taken on the \emph{face} flux and not on
      the cell-centered $\rho v r^2$, because the finite-volume update moves
      the face flux and the centered product is a reconstruction of the
      state;
\item the \emph{chemistry gate}, no cell carrying a composition that is not a
      root of the chemical network, that is no cell in acceptance class 4 or 6
      of Sec.~\ref{sec:ionization:solve}.  Heat and cool in the energy row are
      evaluated on the accepted composition, so on such a cell the row is a
      number computed on a state the chemistry does not describe;
\item the \emph{carrier gate}, where $n(\mathrm{H_2})$ is among the unknowns:
      the carrier row balanced to \texttt{carrier\_resid\_th} of its own
      largest terms.
\end{enumerate}
%
A caller that cannot supply the third makes no statement about the chemistry
and the test rests on the others; the marching stop is such a caller.

\subsection{Certification}
\label{sec:steady:certification}

The four gates above judge the hydrodynamic rows and at most one species row.
A configuration that also carries an elemental transport balance, a level
population or a set of eliminated species would reach its answer with nothing
said about those equations.  \texttt{certification.f90} closes that: it lists
every equation the configuration makes independent and returns, for each, one
of three statuses, \emph{not applicable}, \emph{evaluated} with its residual
measured on \emph{this} state, or \emph{unavailable}.  The third is the point
of the construction: a routine that returns zeros when its frozen background
is not ready is indistinguishable at the call site from a balanced equation,
and an active equation with no evaluator returns nothing at all.  Both now say
so, and a state carrying either is written and not certified.

The inventory covers the three hydrodynamic rows, one balance for each
transported carrier, the ionization stage nucleus sums of H and He, the He/H
partition and the trace elemental transport balances, the
$\mathrm{He}\,2^3S$ and $\mathrm{H}(n{=}2)$ level balances, and the
eliminated-species closure of the coupled cell system.  The stage sums are
the identity $\sum_k F_k(f) + F_{\rm close}(f) = N_{\rm el}(f)$ at every
face, summed over all stages of one element, read from the stationary
evaluation that formed the stage fluxes rather than re-formed, so the number
the report carries is the one the rows were built with; one entry per
element, because the identity is a statement about one element's own nucleus
flux.  It gates, at the floating-point bound of the identity: the identity is
algebraic and has no truncation error, so the whole of what a tolerance on it
can clear is the rounding of the sums, $2(5n_k+2)\epsilon g$ with $g$ the
ratio of magnitude sums measured at the face the measure was taken at, and at
its algebraic ceiling ($3/2$ for one carried stage, $2$ for two) where the
operator reports no $g$.  Constructions that break the identity on purpose
stand ten decades above that bound and the states this code produces at 0.07
to 0.08 of it.  Every row measure is
the \emph{maximum} over physical cells of $|{\rm res}|/\max({\rm scale},
{\rm floor})$, never a volume average: a front or a thin layer must not be
averaged away by the wind that carries most of the volume.  The volume average
is reported beside it and is never decisive.  The measurement has no side
effects; every module array the carrier evaluator writes is held aside around
the call, allocation status included, because a certification is a measurement
of a state and not an operation on it.

\paragraph{The capacity of a report is the size of the inventory.}
A report holds \texttt{cert\_max\_entries} entries, and that number is not a
round constant: it is formed from the counts that set it, three hydrodynamic
rows plus one entry for each transported carrier, two stage nucleus sums, the
He/H partition, one for each trace element, two level balances and one for
each coupled cell system, which is 32 today.  It follows the inventory when a
carrier, a trace element or a cell system is added.  The capacity is exactly
consumed: MEASURED over the regression logs of the binary of record, a
configuration with the ionization stages transported and metals on prints
\texttt{active equations 11 of 32 in the inventory}, so no free slot is left.
The entry writers therefore address only the entry their own call created,
and an entry that does not fit stops the run before any field is written,
naming the equation that did not fit, the capacity and the entry that stands
last and unchanged.  Both halves matter.  What an overfilled report costs is
that a new equation's measure, tolerance and verdict are attached to the name
of the equation before it, and a certification then states a verdict about a
set of equations it does not name; raising the constant would leave that
available one equation later.

The tolerances are one per kind of row, so that the report says which number
refused a state:
$3\times10^{-12}$ for mass, $10^{-8}$ for momentum, $10^{-6}$ for energy,
$10^{-5}$ for an elemental or carrier balance in the wind, and the
composition solve's own root tolerance for the closure and the level
balances.  One number for the three hydrodynamic rows is not possible: on one
converged state the three reach values seven decades apart and their one-ulp
arithmetic floors two decades apart.

Two of these are functions and not constants.  A species row is gated only in
the wind, $r \ge 1.20\,R_p$, and reported without gating below it: the balance
is a cancellation between terms whose ratio is set by the flow, the advective
divergence carrying the row in the wind and the eddy and settling terms
carrying it below the homopause, each of the latter orders of magnitude larger
than their difference, so the precision the same equation can be resolved to is
not one number.  And the mass row's tolerance is a function of the cell,
%
\begin{equation}
  \mathrm{tol}_1(j) = \max\Bigl[3\times10^{-12},\
                         \min\bigl(1,\ c_{\rm round}\,f(j)\bigr)\Bigr],
  \qquad c_{\rm round} = 10,
  \label{eq:masstol}
\end{equation}
%
anchored on the rounding floor of that cell's own continuity row,
%
\begin{equation}
  f(j) = \epsilon_{\rm mach}\,
         \frac{\max_{\rm faces}\bigl[\rho(|v| + c_s)A\bigr]}
              {\max_{\rm faces}\bigl|A F\bigr|}.
  \label{eq:massfloor}
\end{equation}
%
The reasoning is that each interface flux is assembled from two reconstructed
states and from their jumps, and where the column is nearly hydrostatic those
jumps stand many decades below the states, so a correctly rounded interface
flux carries the last bit of the $O(\rho c_s)$ quantities the Riemann solve is
built from and not the last bit of its own value $\rho v$.  The smallest flux
difference the arithmetic can resolve is then
$\epsilon\,\rho(|v|+c_s)A/\Delta V$, which is one ulp in the wind and about
$\epsilon/\mathrm{Mach}$ in a layer whose flux is a tiny residue of the
momentum.  Measured at the base of three reference configurations that floor
spans four decades, $5.7\times10^{-14}$ to $3.2\times10^{-11}$, so a single
fixed number would hold one state to eight ulps of its own base layer and
another to a tenth of one ulp.  The margin $c_{\rm round} = 10$ stands a
factor 3.7 above the largest measured ratio between the step the row takes
when one ulp is added to every density and the estimate \eqref{eq:massfloor}.
The ceiling of unity is not an operative branch on any state measured: the row
measure is the fractional change of the mass flux across a cell, so a state
whose estimated floor reaches 1 is one whose mass row cannot be judged at
all.

\subsection{The Jacobian-free Newton-Krylov solve}
\label{sec:steady:jfnk}

The unknowns of the direct solve are the conserved state on the physical cells,
%
\begin{equation}
  Y = \bigl(\rho_j,\ (\rho v)_j,\ \mathcal{E}_j,\ \{y_{s,j}\}\bigr)_{j=1}^{N},
  \qquad n_{\rm var} = 3 + n_{\rm spec},
  \label{eq:unknowns}
\end{equation}
%
where the optional species rows $y_{s,j}$ are the elemental mass fractions of
the diffusing elements and the transported molecular carriers, registered by
\texttt{set\_transported\_species\_rows} from the configuration.  The ghost
cells are not unknowns: the boundary condition fills them from the interior at
every residual evaluation.  The ionization state and the temperature are
eliminated \emph{locally} inside the residual, that is, given $Y$ the
equilibrium sweep of Sec.~\ref{sec:ionization:solve} re-solves the composition
and returns the consistent heating and cooling of Sec.~\ref{sec:thermal}, so that $F(Y)$ is a residual of size $n_{\rm var}N$ and not of the
full chemical system.

The solve is a Jacobian-free Newton-Krylov iteration in the sense of
\citet{Knoll2004}: the Krylov method needs only the action of the Jacobian on
a vector, which is taken as a forward difference of the residual, so the
nonlocal radiative coupling never has to be assembled.

The outer iteration of \texttt{solve\_steady\_jfnk} is, per step: evaluate the
full residual $F$ at the iterate; build the banded finite-difference Jacobian
of the \emph{frozen-radiation} residual as a preconditioner; solve
$(I/\Delta\tau + J)\,\delta Y = -F$ by right-preconditioned GMRES with the
matrix-vector product taken matrix free, so that the nonlocal radiative
coupling is captured; take a step along a controlled direction; and ramp
$\Delta\tau$.

\paragraph{The banded preconditioner.}
WENO3 reaches two cells on either side and within a cell every variable
couples, so in the flat ordering $Y(n_{\rm var}(j-1)+k)$ a column couples to
rows within $|{\rm row} - {\rm col}| \le 3n_{\rm var}-1$; the half-bandwidths
are that number and a graph coloring with stride
$n_{\rm color} = k_l + k_u + 1 = 6n_{\rm var}-1$ makes the row supports of
same-color columns disjoint, so that the frozen-radiation residual, which is
strictly local, is captured exactly by that many finite-difference probes.  For
$n_{\rm var} = 3$ this is 8 and 17.  The band is factored by LAPACK
\texttt{dgbtrf} and applied as the right preconditioner of the Krylov cycle.
The finite-difference and Krylov probe steps are taken in the diagonally
scaled variables of \texttt{cell\_state\_scales}, the characteristic magnitude
of each unknown in its own cell; a probe point that leaves the set of states
the code can describe, a chemical balance row going to \texttt{NaN} or
overflowing, is answered by halving the step and sampling closer to the
iterate, at most six times, and never by accepting a non-finite residual or
inventing a derivative.

\paragraph{The step control.}
Where the system carries no species row the step is the Krylov direction with
a backtracking line search on the merit $\|F/D\|_2$, with a Grippo
non-monotone reference, the worst true merit of the last five iterates, so
that sideways steps are admissible and the solve can leave a bad
neighborhood.  Where the system carries a species row the step control is a
scaled trust region with a Powell dogleg: the path runs from the origin to the
Cauchy point $s_U$, the minimizer of the linear model along the steepest
descent direction, and from there to the Krylov point $s_N$, and the step is
where that path first meets the ball of radius $\Delta$,
%
\begin{equation}
  s = a_U s_U + a_N s_N,
  \qquad
  \begin{cases}
    a_U = 0,\ a_N = 1                 & \|s_N\| \le \Delta,\\
    a_U = \Delta/\|s_U\|,\ a_N = 0    & \|s_U\| \ge \Delta,\\
    a_U = 1-\tau,\ a_N = \tau         & \text{otherwise},
  \end{cases}
  \label{eq:dogleg}
\end{equation}
%
with $\tau$ the root of $\|s_U + \tau(s_N - s_U)\| = \Delta$ in $[0,1]$.
Returning the two coefficients rather than the vector is deliberate: the image
$As$ follows from $As_U$ and $As_N$ by linearity and costs no further
matrix-free product.  The Krylov leg is solved without the pseudo-transient
shift, $As_N = -r_0$ rather than $(I/\Delta\tau + A)s_N = -r_0$, because the
dogleg falls monotonically along its path only when that leg is the model's own
minimizer.  A species unknown also carries a box, the fraction-to-the-boundary
rule holding the step inside the physical range of a mass fraction, and an
element with no reservoir is excluded from the registry altogether: its
unknown would sit exactly on its own lower bound with no room below, and the
boundary rule then returns zero for any direction with a component pointing
there, so the whole Krylov direction would be refused.

\paragraph{The ledger and the stop.}
The best iterate reached so far is kept, with its composition, its flux spread
and whether it met the gates.  The non-monotone search accepts sideways steps,
which is what lets the solve leave a bad neighborhood and also what lets it
circle, so when the best iterate has not improved for 20 outer iterations the
solve returns to it and switches to a monotone Armijo acceptance, and if that
buys nothing in another 20 it stops and reports that it stagnated, a different
statement from ``no descent direction exists''.  The loop-top stop requires
the certified rows of Sec.~\ref{sec:steady:certification} and not only the
merit: the two are different functionals, and a window on the merit alone
bounds nothing about the quantity the state is judged by.  The return flag is
0 when the returned state met the acceptance gate and every judged row, 1 at
the iteration limit, 2 when the solve ended without a state meeting the gate,
and 3 at an iterate where no feasible descent direction exists.  A flag of 2
is not a solver failure: the solve converged and the certification refused the
state it hands back, naming the row, and that state is written with the
refusal recorded rather than marched on.

\subsection{Pseudo-transient continuation}
\label{sec:steady:ptc}

\texttt{solve\_steady\_ptc} is the same residual solved with the banded
Jacobian used as the \emph{operator} and not only as a preconditioner: per
iteration $(I/\Delta\tau + J)\delta Y = -F$ is solved by LAPACK banded LU,
the step is taken with a backtracking line search and a positivity test, and
$\Delta\tau$ follows the switched-evolution-relaxation ramp
$\Delta\tau \leftarrow \Delta\tau\,\|R\|_{\rm old}/\|R\|_{\rm new}$, cut on a
failed step.  As $\Delta\tau \to \infty$ this is Newton's method; at small
$\Delta\tau$ it behaves like explicit relaxation, which is the robust startup
the continuation exists for.  The route is not available with a carrier
unknown: it builds its Jacobian from the frozen residual, which by
construction performs no equilibrium sweep and therefore has no carrier row,
and giving it one would mean assembling that row from a background nothing in
the routine refreshes.  The caller is told so rather than being handed a
silently three-row answer.

The starting $\Delta\tau$ matters: the marching handoff starts the
continuation from $\Delta\tau_0 = 1$ and the two direct routes from the CFL
step.

\subsection{The partitioned outer iteration}
\label{sec:steady:partitioned}

A species the run moves by an operator-split transport step instead of by a
Newton row, which is the diffusing elements always and the molecular carriers
whenever \texttt{Coupled carrier solve} is off, has a stationary balance no
hydrodynamic solve touches; a single steady solve would freeze it at whatever
state it was handed.  \texttt{steady\_wind\_with\_element\_diffusion} in
\texttt{EXHALE\_main.f90} alternates the two: solve the wind at fixed
composition, relax the composition to its own steady state in that wind,
solve again. ``Fixed composition'' means the \emph{transported} part of it:
the element fractions the diffusion operator moves and the carrier densities.
The local part, the ionization state of every cell and the molecular
reaction balance at the carriers' densities, is re-solved inside every
evaluation of the hydrodynamic residual (Sec.~\ref{sec:steady:residual}),
so the wind solve carries its own chemistry and only the transport of
composition is alternated. A run with neither element diffusion nor carrier
transport therefore has one outer pass: the wind solve alone.  The composition relaxations are the fixed-wind element
relaxation of \texttt{binary\_element\_diffusion.f90}
(Sec.~\ref{sec:mol:diffusion}), under-relaxed with a factor $\omega$ starting
at $0.5$, and the fixed-wind carrier relaxation of
\texttt{diffusive\_photochemistry.f90} (Sec.~\ref{sec:mol:carrier}), bounded
by a movement trust rather than damped.  The advecting flux handed to the element operator is the face
mass flux of the state just solved, read from the row terms that state
assembled, so the operator is given its wind and does not reach into the
residual for one.

What is accepted is only a state on which the certification of the
\emph{full} set of active equations passes, evaluated on the refreshed state
the routine would hand back.  A hydrodynamic solve returning 0 while an
alternated species row refuses is not an accepted state: that flag is a
statement about the rows in the Newton registry, and the alternated rows
stand outside it by construction.  Conversely a nonzero hydrodynamic flag does
not end the loop by itself, because the composition update that follows it
still makes progress.

\paragraph{The progress control.}
A pass makes progress when the joint distance of the state to the
certification falls, or, where that distance is carried by a row that is not
hydrodynamic, when the composition's own residual falls; three consecutive
passes without progress end the loop rather than spend the budget (20 passes,
\texttt{EXHALE\_OUTER\_PASSES}) on a fixed point the alternation is not
approaching.  The endings are named: the joint certification passes; a row of
the state is not finite; the element composition update is refused; the
joint measure stops falling; the pass budget ends uncertified.

The two step lengths of the composition half, the element under-relaxation
$\omega$ (start $0.5$, \texttt{EXHALE\_DIFF\_OMEGA}; floor $0.125$) and the
carrier movement bound, the largest relative change of one cell's particle
count $n_{\rm tot}+n_e$ one pass may make (start $10^{-2}$,
\texttt{EXHALE\_CARRIER\_TRUST}; floor $10^{-3}$; ceiling $0.5$), answer to
the hydrodynamic solve that consumes the update (revised 2026-09-27).
\emph{A pass whose carrier update ended on its movement bound is not a pass
without progress} as long as the solve that consumed it converged and the
joint distance rose by no more than a factor $1.1$: the state moved by the
whole bound, and the relative measures cannot show it because residual and
scale fall together.  MEASURED on the LHS~1140\,b He/H\,$=1.8$ state with
transported ionization stages, every pass ended on the bound while the gated
measure stood at $2.17$ to $2.20\times10^{-2}$; the earlier control read that
as stagnation, halved the bound to its floor and ended the loop, while the
same state with the bound held at $0.5$ certified in 14 passes.  Such a pass
is neither counted toward the ending nor allowed to shorten a step, and the
bound doubles, up to $0.5$.  The bound halves when a bound-limited update
raised the joint distance of a solved wind by more than $1.1$, when it
preceded the first failed solve after a converged one, and when a pass
without progress moved the carriers short of the bound on a solved wind; a
relaxation that kept no step leaves it where it is.  $\omega$ halves on a pass
without progress whose wind was solved and on the first failed solve after a
converged one, and doubles back toward its starting value after a converged
solve on a pass with progress.  A failed solve that follows another failed
solve did not start from a solved wind, so the update is not what failed it
and neither step is shortened.  \texttt{EXHALE\_CARRIER\_TRUST\_HOLD} freezes
the bound for every pass.  Recorded as open: the starting bound for
transported stages (the state above certifies from $0.5$ and not from
$10^{-2}$), and the molecular He/H\,$=0.083$ continuation, where a growing
bound makes the H$_2$ row of the outermost cells overshoot and neither control
certifies.

\paragraph{The state of every pass, kept.}
Every pass after which the iteration goes on writes the state the next pass
starts from, rebuilt from $(u, f_{\rm sp})$ in the order of the checkpoint
restoration with every side effect on the caloric mixture undone, to
\texttt{output/Hydro\_ioniz\_last\_pass.txt} and
\texttt{output/Ion\_species\_last\_pass.txt} by the writers of the final
state, so a run stopped from outside keeps its last completed pass.  The pair
carries \texttt{certified=F cert\_reason=pass\_snapshot\_p<n>} in both halves,
because the certification of a pass is taken before its composition update
and the snapshot is the state after it, and a \texttt{\# pass\_snapshot} line
with the movement bound and $\omega$ of the next pass, which a restart does
not inherit.  Each half is written under a temporary name and renamed, and a
pair mixed from two passes states two reasons and is refused by the loader.
Nothing the run reads afterwards changes: the final state files of a run are
bitwise those of the binary before the snapshot existed.

Two rules keep the state self-consistent.  The progress control stands
\emph{ahead} of the composition update, so the pass that ends the iteration
takes no update and what is handed back is the very state the certification of
that pass was evaluated on; and the last pass of the budget takes no update
either, an update existing only to be consumed by the next hydrodynamic solve.
Taking it would replace a stationary wind by a state whose hydrodynamic rows
the update itself spoiled, by eleven decades in the mass row on a molecular
configuration.

With \texttt{Coupled carrier solve: True} the alternation does not arise for
the carriers at all: $n(\mathrm{H_2})$ becomes a fourth Newton unknown per
cell, one solve returns a wind and a carrier partition that are steady states
of each other, and the outer loop runs once.  A third value,
\texttt{On stall} (2026-09-17), runs the alternation and hands the state to
that block at the pass where the alternation gives up: the joint distance has
not fallen in three consecutive passes AND the carrier relaxation of this pass
and of the two before it ended on the movement bound.  What is known about the
two routes where they have been compared: on a certified molecular state the
alternation certifies at outer pass 12 while the block moves its residual by
0.6 per cent in the same wall clock, every one of its linear solves ending with
the Krylov subspace exhausted, so the block as it stands is limited by its
preconditioner and not by the species box.

\paragraph{What the alternation cannot do, measured.}
The movement bound on the carrier relaxation is not what holds the refusing
molecular states: relocating the cells that control it only moves the control
to the edge of whatever region is exempted, and the refusing row is worse
afterwards.  What the measurements point at is the splitting itself. The bound
exists because the wind is held while the carriers relax, and a composition
good enough for the carrier rows is already outside the range in which holding
the wind is admissible: the one relaxation that reached the carriers' own fixed
point moved the base particle count almost seven times the bound, and the
hydrodynamic solve that had to follow it did not converge, its mass rows
outside tolerance in every cell over a hundred iterations.  With no alternated species at
all, which is what an atomic run and a fully coupled run both give, the body
is one steady solve followed by the state refresh and the flag is the solve's
own.
